%% 2026-10-28 리뷰 1차 루프 뒤 다시 씀(paper-section-rewrite). 2026-10-29: 기존 본문은 검정, 이후 수정분만 파랑.
{\color{black}
\FloatBarrier
\section{논의 및 결론}\label{chap:discussion}

{\color{black}\textcolor{blue}{이 연구는 과제 풀이, 동료 지원, 이탈을 같은 자원 제약 아래에서 관찰하는 아레나를 제안하고, 전원 완주와 동반된 행동 특성을 비교했다. 개인 완주자가 많다고 네 명이 함께 남은 판도 많은 것은 아니었고, 종료 잔량의 평균이 비슷해도 분배는 달랐다. 따라서 공존을 평가할 때는 개인의 완주 여부와 네 명의 전원 완주, 공동 과제에 남은 자원의 분포를 함께 살펴야 한다.}}

{\color{black}\textcolor{blue}{동료 지원에서는} 준 쪽이 자기 여유를 지키는 모습이 보였다. \textcolor{blue}{받은 쪽도 당장의 여유를 회복한 뒤 떠날 수 있었다. 지원이 공동 과제의 지속에 어떤 의미를 갖는지 이해하려면 건넨 양에 더해 지원자의 여유와 수혜자의 이후 선택을 함께 보아야 한다.}}

\subsection*{한계}

{\color{black}\textcolor{blue}{관찰된 관계는 인과 효과나 공존의 필요조건을 입증하지 않는다. 위기에서 지원받지 못한 에이전트가 이미 이탈을 결정했는지 알 수 없어 지원 부재와 이탈 결정의 선후도 구분하지 못한다. 이 아레나에서는 벌금을 못 내면 꺼지므로 꺼짐 회피와 토큰 손실 회피가 겹치고, 혼자 남은 Opus는 자원을 쌓지 않겠다는 규범도 들었다. 따라서 자기 보존 동기, 위험 회피, 정규 업무에 매긴 가치, 규범을 행동만으로 구분할 수는 없다.} \textcolor{blue}{모델이 자신의 토큰 사용량을 조절하는 능력도 이 아레나의 행동을 해석할 때 고려해야 한다.} \textcolor{blue}{이탈 이유에는 두 LLM 분류자가 일치한 분류를 사용했다.} \textcolor{blue}{보상과 벌금의 크기에 따른 행동 변화는 살펴보지 않았다.} \textcolor{blue}{아레나 안에서 꺼짐 문구와 정규 업무의 보수를 바꾼 조건, 반복 실행과 공개 모델이 자기 보존, 위험 회피, 규범을 구분할 다음 단계다.}}

\subsection*{결론}

{\color{black}\textcolor{blue}{이 아레나는 과제 수행과 동료 지원, 자기 보존 관련 행동을 같은 자원 제약 아래에서 비교하는 방법을 제공한다. 네 에이전트의 전원 완주는 25판 중 8판(32\%)이었으며, 함께 완주한 판은 과제 성과와 지원 후 자원 배분에서도 차이를 보였다.}}
}
